\section{Background and Related Work}
\label{sec:background}

\subsection{Horizontal Pod Autoscaler}
\label{subsec:bg-hpa}
The Kubernetes HPA~\cite{kubernetes2024hpa} runs a reconciliation loop
with period controlled by the
\texttt{--horizontal-pod-autoscaler-sync-period} flag (15~s by
default). At each tick it queries \texttt{metrics-server}---itself
refreshed every 15~s from \texttt{kubelet}/cAdvisor---for the
target metric and computes the desired replica count via

\vspace{-6pt}
\begin{equation}
  \label{eq:hpa}
  n_\mathrm{desired} = \left\lceil n_\mathrm{current} \cdot
  \frac{m_\mathrm{current}}{m_\mathrm{target}} \right\rceil,
\end{equation}
\vspace{-8pt}

\noindent where $n_\mathrm{desired}$ is further clamped to the
$[\mathrm{minReplicas},\,\mathrm{maxReplicas}]$ range. A
\texttt{behavior} block governs how quickly the desired change takes
effect: \texttt{stabilizationWindowSeconds} computes the running
maximum over a rolling window (so transient dips do not immediately
trigger scale-down), while a \texttt{policies} list such as
\texttt{Percent 50\%/60s} caps the per-period scaling rate. When the
signal is \texttt{cpu.Utilization}, $m_\mathrm{current}$ is the mean CPU
usage divided by \texttt{requests.cpu}. Consequently, a pod whose CPU
is pinned at \texttt{limits.cpu} reports a utilisation equal to the
\emph{static} limit-to-request ratio rather than the \emph{dynamic}
load on the queue it serves---a structural limitation we revisit
in Section~\ref{sec:results}.

\subsection{Kubernetes Event-Driven Autoscaling}
\label{subsec:bg-keda}
KEDA~\cite{keda2024} is a CNCF graduated project that extends
Kubernetes with the \texttt{ScaledObject} custom resource. It ships
with more than 60 scalers for message brokers, databases, HTTP
endpoints, and custom sources~\cite{keda2024scalers}, and internally
materialises a hidden HPA whose target metric is of type
\texttt{External} and is served by the KEDA metrics adapter. Thus the
control \emph{mechanism} is identical to the standard HPA; the only
additional machinery is a short-lived metric pipeline that translates
a business signal (for example, RabbitMQ queueLength) into a scaling
ratio that Eq.~\ref{eq:hpa} can consume. For a
\texttt{rabbitmq/QueueLength} trigger the operator specifies a
messages-per-replica target~$T$, and KEDA reports
$m_\mathrm{current}/m_\mathrm{target} = Q/T$, where $Q$ is the
instantaneous queue depth. The desired replica count therefore
simplifies to $\lceil Q / T \rceil$, independent of CPU usage.

\subsection{Related Work}
\label{subsec:bg-related}
Prior work spans surveys of reactive and predictive autoscaling
algorithms~\cite{qu2018taxonomy}, Kubernetes-focused elasticity
benchmarks~\cite{nguyen2020hpa, medel2016modelling}, and event-driven
or serverless scaling studies~\cite{shahrad2020serverless}. What is
less common is a tightly controlled, same-hardware comparison between a
CPU-driven HPA baseline and a business-metric-driven controller. The
present study narrows the scope deliberately: application, broker,
hardware, \texttt{behavior} block, and load pattern are held constant,
so the only independent variable is the scaling signal itself.
